\documentclass[10pt,a4paper]{amsart}
\usepackage[margin=19mm]{geometry}
\usepackage{amsmath,amssymb,mathtools}
\usepackage[scr=rsfs,cal=euler]{mathalfa}
\usepackage{xfrac}
\usepackage[unicode,hidelinks,bookmarks=false]{hyperref}
\newcommand{\ie}{i.e.\ }
\newcommand{\cf}{cf.\ }
\newcommand{\resp}{respectively\ }
\newcommand{\Subsection}{\subsection}
\providecommand{\nobreakdash}{\nobreak\hbox{-}}
\setlength{\emergencystretch}{2em}
\multlinegap0pt
\usepackage{bm}
\usepackage{bbm}
\usepackage{dsfont}
\usepackage{xspace}
\usepackage{cancel}
\usepackage{mathtools}
\usepackage{amsthm}
\makeatletter
\@ifundefined{theorem}{\newtheorem{theorem}{Theorem}}{}
\makeatother
\newcommand{\dnint}[1]{\left\|{#1}\right\|}
\renewcommand{\mod}{\;\mathrm{mod}\;}
\title[A Note on Diophantine Approximation with Restricted Denomina]{A Note on Diophantine Approximation with Restricted Denominators}
\author{Chance Sanford}
\date{Source: arXiv:2606.02620v1 (CC BY 4.0)}
\begin{document}
\maketitle

\noindent\emph{Source and licence.} A Note on Diophantine Approximation with Restricted Denominators. Version arXiv:2606.02620v1 (CC BY 4.0), distributed under the Creative Commons Attribution 4.0 International licence, as recorded in the arXiv licence metadata for this exact version and shown on the abstract page https://arxiv.org/abs/2606.02620v1. Full rights notice: licenses/arxiv-2606.02620-CC-BY-4.0.txt.\par\smallskip

\noindent\textit{Editorial scope.} Counted unit: the Theorem whose source label is \texttt{thm:set\_diff\_ddense} in the article (source0.tex lines 154--156, source label \texttt{thm:set\_diff\_ddense}), delivered together with the article's own complete original proof of it (source0.tex lines 157--215, one \texttt{proof} environment of 59 lines). No mathematics is written, repaired or re-derived: statement and proof are the article's own text line for line. Required context retained uncounted: none. Every definition, display and label of those lines is retained unchanged. Omitted from this package and not redistributed: 1--153, 216--284. Nothing inside a retained excerpt is omitted or altered: every retained body is delivered exactly as the article's lines are written, so no omission receipt of any kind accompanies this package. Citations: the retained excerpts carry no citation command at all, so no bibliography entry of the article is transcribed and no reference is redistributed. Reference adaptation: none was needed and none was made; every reference inside the delivered unit resolves inside it, so no unresolved reference or citation is produced. Printed slips kept literal: no printed word, symbol, hypothesis or number of the counted statement or of its proof was altered. Layout and class: the article's own class is \texttt{amsart} with packages including inputenc, amsmath, amsfonts, amssymb, amsthm, mathrsfs, geometry; the wrapper is the lane's pinned \texttt{amsart} editorial wrapper at 10pt on A4, which restates the theorem-like environments the retained text needs and adds the article's own macro definitions that the retained text needs (\texttt{\textbackslash dnint}, \texttt{\textbackslash mod}), with extra wrapper packages (bm, bbm, dsfont, xspace, cancel, mathtools). The delivered numbering is the unit's own. Assets: no retained span contains a figure environment, a graphics-inclusion command, a PDF-page inclusion, a table or a TikZ picture; no image file is copied into this package, the package declares no asset dependency and no image file is redistributed with it. Licence: version arXiv:2606.02620v1 of this article is distributed under the Creative Commons Attribution 4.0 International licence, as recorded in the arXiv licence metadata for this exact version and shown on the abstract page https://arxiv.org/abs/2606.02620v1. A full rights notice is delivered at licenses/arxiv-2606.02620-CC-BY-4.0.txt. Characters: every retained excerpt is pure ASCII, so the delivered engine, XeLaTeX with its default Latin Modern text font, reports no missing character.
\begin{theorem}\label{thm:set_diff_ddense}
Let $\mathcal{A} \subseteq \mathbb{N}$ and $\delta \in [0,1]$. If $\mathcal{A}(n) \ll n^{\delta}$ for sufficiently large $n$, then the set $\mathbb{N} \,\backslash \,\mathcal{A}$ has a Diophantine density of $1-\delta$.
\end{theorem}

\begin{proof}
Fix a coprime $(p,q) \in \mathbb{Z} \times \mathbb{N}$ with $q$ sufficiently large.  Our first task is to analyze the multiset
\begin{equation*}
    \Psi(p,q):= \left\{\dnint{\frac{k p}{q}}\,:\,  1 \leq k \leq q \right\}.
\end{equation*}
Once we have done that, it will be straightforward to bound the minimum element in the submultiset
\begin{equation*}
    \left\{\dnint{\frac{k p}{q}}\,:\,  1 \leq k \leq q, k \in \mathbb{N}\,\backslash\,\mathcal{A} \right\}.
\end{equation*}
To this end, note that for each $k$ we have
\begin{equation*}
\dnint{\frac{kp}{q}} = \frac{\left\vert{kp-N_{k}q}\right\vert}{q},
\end{equation*}
where $N_{k}$ is the nearest integer to $kp/q$.  Since $\dnint{x} \leq 1/2$ for any real number $x$, we obtain the bound
\begin{equation*}
\vert{k p -  N_{k} q }\vert \leq \frac{q}{2}.
\end{equation*}
Consequently, we find that
\begin{equation*}
\dnint{\frac{kp}{q}} \in \left\{0,\frac{1}{q},\frac{2}{q},\dots,\frac{\lfloor{q/2}\rfloor}{q}\right\}.
\end{equation*}
Next, by definition we have
\begin{equation*}
    \dnint{\frac{p}{q}} = \dnint{\frac{p \mod q}{q}} = \frac{1}{q}\min\left\{p \mod q, q - (p \mod q)\right\}.
\end{equation*}
Thus, for $p_{1},p_{2} \in \mathbb{Z}$, we have
\begin{equation*}
    \dnint{\frac{p_{1}}{q}} = \dnint{\frac{p_{2}}{q}},
\end{equation*}
if and only if $p_{1} \equiv \pm p_{2} \mod q$.

For our purposes, this enables us to determine when
\begin{equation*}
    \dnint{\frac{k_{1} p}{q}} = \dnint{\frac{k_{2} p}{q}},
\end{equation*}
for $1\leq k_{1},k_{2}\leq q$.  Namely, since $p$ is invertible mod $q$, the equality holds if and only if $k_{1} \equiv \pm k_{2} \mod q$.  Or equivalently, since $k_{1},k_{2} \leq q$, if and only if $k_{1} = k_{2}$ or $k_{1} = q-k_{2}$.

From this fact we deduce that each number in the set
\begin{equation*}
\left\{0,\frac{1}{q},\frac{2}{q},\dots,\frac{\lfloor{q/2}\rfloor}{q}\right\}
\end{equation*}
occurs at least once and at most twice in the multiset $\Psi(p,q)$.  In particular, $0$ occurs once for every $q$ and $1/2$ occurs once for even $q$, with all other elements occurring exactly twice.

Now that we understand the structure of the multiset $\Psi(p,q)$, our next goal is to determine what can be said about the minimum remaining element after removing the submultiset of elements which correspond to elements of $\mathcal{A}$ less than or equal to $q$.  To accomplish this, recall that the requirements of the theorem dictate that for large enough $q$ and $c > 0$,
\begin{equation*}
\mathcal{A}(q) \leq c q^{\delta}.
\end{equation*}
Suppose that we remove the smallest $\lfloor{c q^{\delta}}\rfloor$ elements from the multiset $\Psi(p,q)$.  Since each number $k/q$ occurs at least once for $0 \leq k \leq \lfloor{q/2}\rfloor$, the smallest remaining element may be roughly bounded above by ${\lfloor{c q^{\delta}}\rfloor/q}$.

In other words, we find that
\begin{align*}
\min_{v\in\mathbb{N} \,\backslash \,\mathcal{A}, v\leq q}\dnint{\frac{v p}{q}} \leq \frac{\lfloor{c q^{\delta}}\rfloor}{q} \leq \frac{c}{q^{1-\delta}}.
\end{align*}
We have thus demonstrated that there exists a constant $C > 0$, depending only on $\mathbb{N} \,\backslash \,\mathcal{A}$, such  that for any coprime $(p,q)\in\mathbb{Z}\times\mathbb{N}$ with $q$ sufficiently large,
\begin{equation*}
    \min_{v\in\mathbb{N} \,\backslash \,\mathcal{A}, v\leq q}\dnint{\frac{v p}{q}} \leq \frac{C}{q^{1-\delta}}.
\end{equation*}
This completes the proof.
\end{proof}

\end{document}
